\section{Using two later-row bounds simultaneously}

\subsection{O-018.a: a joint-row gate}
In the BIP numerical range $n\geq m^{25}$, every actual
multiplicity obstruction satisfies
\[
                         \lambda_3\geq139
                           \quad\hbox{or}\quad\lambda_{52}\geq3.
                                                               \tag{O-018.1}
\]
The numerical constant is a sufficient bound obtained from
the count below; it is not claimed optimal.
A-025 subsequently strengthens $139$ to $557$ by using
two support trees and retaining the total-degree constraint.
The elementary count here remains useful for the explicit
full-profile example below.

\begin{proof}
The earlier reductions already give $m\geq190$ and
$52\leq r=\ell(\lambda)\leq m^2$ for an actual obstruction.
Suppose instead that $\lambda_3\leq138$ and
$\lambda_{52}\leq2$. In a projected degree-$m$ form,
the nested support F-016 relevant to a highest-weight
function has total degree at most $138$ in variables
$3,\ldots,r$, and at most two in variables $52,\ldots,r$.
There is one earlier variable, $x_2$, forty-nine middle
variables, and $r-51$ last variables.

Let $\Gamma$ be the set of all nonfirst-variable multi-indices
of degree at most $m$ satisfying the full nested bounds.
It is nonempty and downward closed. Ignoring coupled degree
conditions only increases its size, so
\[
 |\Gamma|\leq(m+1)\binom{187}{49}\binom{r-49}{2}
                           \leq\binom{187}{49}m^5.     \tag{O-018.2}
\]
The second inequality uses $m+1\leq2m$ and
$\binom{r-49}{2}\leq r^2/2\leq m^4/2$.
Exact integer arithmetic gives
\[
                         \binom{187}{49}\leq190^{20}.
\]
Since $m\geq190$, (O-018.2) is at most $m^{25}\leq n$.

Place all of $\Gamma$ in one monomial tree, and take the
other tree to consist of its root. O-015.a gives a determinant
of size $|\Gamma|$ for every retained padded form.
For every determinant-ideal highest-weight vector, first align
the padding with $x_1$ as in M-009, then delete the coefficient
monomials outside $\Gamma$. F-016 preserves its evaluation,
and the retained form is the determinant just constructed.
Thus every such vector vanishes on the permanent closure.
The inclusion of all highest-weight ideal multiplicity spaces
gives $q\leq K_n$, a contradiction. This proves (O-018.1).
This is also a special case of the general support-factorization
criterion F-018, with one factor equal to the whole support
and the other consisting only of zero.
\end{proof}

\subsection{O-018.b: a four-copy example beyond every single-cut containment test}
Set
\[
 m=2^{200},\quad r=m^2,\quad n=m^{25}=2^{5000},\quad d=3n,
\]
and define
\[
 \lambda=(3n^2-2n-2r-1172,\ 2n,\ 26^{49},\ 2^{r-51}).
                                                               \tag{O-018.3}
\]
It has $r$ rows and size $nd$. Then
\[
                              4\leq q^m_\lambda(d[n])
                                              \leq K_n(\lambda).
                                                               \tag{O-018.4}
\]
The exact values of these multiplicities are not asserted.
Every single-cut universal coefficient family associated to
this partition has dimension greater than the corresponding
determinant-pencil closure. Thus the comparison in (O-018.4)
uses information beyond every one-cut containment test, not
only beyond its uniform staircase corollaries.

\paragraph{The complete support fits.}
For this partition the full nested degree-$m$ support says
that the middle forty-nine variables and the last $r-51$
variables together have degree at most $26$, while the
last variables alone have degree at most two. The exponent
of $x_2$ is limited only by the total degree $m$, because
$\lambda_2=2n>m$.
The exact number of retained indices is
\[
 G=\sum_{t=0}^2\binom{r-52+t}{t}
       \sum_{j=0}^{26-t}\binom{48+j}{j}(m-j-t+1).
                                                               \tag{O-018.5}
\]
Here $t$ is the last-variable degree, $j$ the middle degree,
and the final factor counts powers of $x_2$.
In particular,
\[
 G\leq(m+1)\binom{75}{26}\binom{r-49}{2}
          \leq\binom{75}{26}m^5<2^{75}m^5<m^{25}=n.
\]
The one-tree construction in O-018.a therefore proves
$q\leq K_n$ for this partition as well.

\paragraph{At least four actual permanent copies.}
Take the four independent highest-weight classes on the
permanent closure of degree six and weight $(6n-6,6)$,
supplied by M-012.2. Multiply all four by the common factor
\[
                   W_2^{n-16}W_{51}^{12}W_r
                                       F(e_1)^{n-r-586}.        \tag{O-018.6}
\]
All exponents are nonnegative. Its degree is
$2n-32+612+r+n-r-586=3n-6=d-6$.
After adding the seed weight, its second row is
$6+2(n-16)+24+2=2n$; rows three through fifty-one
are $24+2=26$; the remaining nonfirst rows are two.
The first row is determined by the total degree, giving
exactly (O-018.3).

O-011 supplies a nondegenerate permanent Hessian at a point
where the value of the form and its radial Hessian entry
are nonzero. Extend that radial direction to an orthogonal
basis. Every leading Hessian minor, including $W_2,W_{51},W_r$,
is then nonzero. Thus (O-018.6) is a nonzero class on the
permanent closure. Its reduced coordinate ring is a domain,
so multiplication by this one class preserves the independence
of all four seed classes. This proves the lower bound in
(O-018.4); it does not require all seed classes to be nonzero
at the particular Hessian point.

\paragraph{Every single cut is too large as a universal family.}
For $2\leq k\leq r$, let $\mathcal U_k$ be the whole
space of degree-$m$ forms whose degree in variables
$k,\ldots,r$ is at most $\min(m,\lambda_k)$, with no
other later-row restriction. Also allow $k=r+1$, with an
empty last group and no restriction. Its padded copy has
the same dimension. We show
\[
                         \dim\mathcal U_k>m^{52}
                                    >(r-2)n^2+2         \tag{O-018.7}
\]
for every cut. By O-010, none of these whole padded families
can be contained in $\Omega_n$, even allowing limits.

If $k=2$, all degrees through $m$ are allowed in the last
group. If $k>r/2$, at least $r/2-2$ early variables are
unrestricted. In either case squarefree monomials of degree
twenty-seven in an available group, with the remaining degree
on $L$, give
\[
 \dim\mathcal U_k\geq\binom{r/2-2}{27}
       \geq\frac{(r/4)^{27}}{27!}
       =\frac{m^{54}}{2^{54}27!}>m^{52}.                \tag{O-018.8}
\]
The factorial inequality follows from
$2^{54}27!<2^{54+135}<2^{400}=m^2$.

For $3\leq k\leq51$, the last-variable cap is twenty-six
and at least one early variable is available. Choose
twenty-six distinct last variables and any power through
$m-26$ of one early variable. This gives
\[
 \dim\mathcal U_k\geq(m-25)\binom{r-50}{26}
       \geq\frac{m^{53}}{2^{27}26!}>m^{52},            \tag{O-018.9}
\]
using $m-25\geq m/2$, $r-75\geq r/2$, and
$2^{27}26!<2^{157}<m$.

For $52\leq k\leq r/2$, use fifty early variables and
last-variable degree exactly two. The dimension is at least
\[
 \binom{m+48}{50}\binom{r/2+1}{2}
       \geq\frac{m^{54}}{2^{53}50!}>m^{52},            \tag{O-018.10}
\]
because $2^{53}50!<2^{353}<m^2$.
These ranges exhaust all cuts, including the empty-tail
case in (O-018.8). Finally,
$(r-2)n^2+2=(m^2-2)m^{50}+2<m^{52}$.
This proves (O-018.7).

\paragraph{Other explicit gates and scope.}
The BIP threshold holds with equality. All F-017 staircase
inequalities hold: the rows through fifty-one in question
are twenty-six and row fifty-two is two. The full-multiplicity
stabilization criteria requiring $n\geq\lambda_2$ fail
since $\lambda_2=2n$. The universal quadratic and grouped
Chow degree criteria fail: their exact thresholds at $L=dm$
are at least $\binom d2>n$.
Grenet's full containment criterion fails because
$m=2^{200}>5001$ and hence $2^m-1>2^{5000}=n$.
The older coordinate-box test fails as well: some part in
every split has at least $r/2$ variables, so its count is
at least $\binom{r/2}{27}>n$.
The coarse padding inequalities hold since
$|\bar\lambda|=2n+2r+1172\leq md$ and $r=m^2$.

The finite certificate checks the exact count (O-018.5),
the three uniform lower bounds covering every cut, the
joint-row constant in (O-018.2), and every stated numerical
comparison. No enumeration of the enormous number of cuts
is needed. This is an excluded shape with actual multiplicity
at least four, and a strict reduction of the earlier explicit
conditions. It is not an actual $q>K_n$ obstruction.

\paragraph{Geometric dictionary.}
The entire nested profile selects a downward-closed finite
torus-weight coefficient subscheme. Passing from individual
row caps to this profile is a calculation with exterior
matrix coefficients and their Cauchy span, as in F-016;
the exact background locators are Macdonald I(4.3), p.~63,
and I Appendix~A(5.3)--(5.6), pp.~155--156, and SGA~3 I,
Proposition~4.7.3, pp.~35--36. The one-tree substitution is
an integral morphism to the linear-matrix scheme. The actual
comparison is obtained only after retaining every
highest-weight ideal copy in characteristic zero. The
dimension obstruction uses the image closure, while the
four permanent copies use multiplication in its integral
coordinate ring; these are distinct arguments on the actual
schemes, not deductions from a common virtual $K_0$ class.
